\section{Smoothed deletion measures and their local comparison}
\label{palm:section}\label{sec:local-palm}

In this section $\Phi$ is a normalized nonnegative ground function.
For a cell $v$ in either $\mathcal C_d$ or $\mathcal C_R$, define
\begin{equation}\label{palm:rms-density}
 s_v(Y')=\left(\frac V{|v|}\int_v\Phi(y,Y')^2\,\dd y\right)^{1/2},
 \qquad P_v(\dd Y')=s_v(Y')^2\,\dd Y'.
\end{equation}
The measure $P_v$ describes the remaining particles after deletion of
a particle in $v$, with its position averaged over the cell.  Its mass
need not be one.  We shall prove the following estimate for neighboring
small cells; Section~\ref{asm:section} uses its fourth power to compare
cell averages of $\Phi$ with their root mean squares.

\begin{proposition}[Local comparison of smoothed deletion measures]
\label{palm:local-comparison}
With the exponents and partitions of Proposition~\ref{en:energy-outputs},
every normalized nonnegative bosonic ground function satisfies
\begin{equation}\label{palm:local-rms-comparison}
 \frac1{|\mathcal C_d|}\sum_{v\sim w}
       \int\frac{|s_v-s_w|^4}{s_v^2+s_w^2}\,\dd Y'
                 \le C_jD^{-1-u/8}.
\end{equation}
Here $v\sim w$ ranges over ordered adjacent $d$-cells, and the
quotient is zero when both functions vanish.  The sum may be restricted
to pairs in the same $R$-cell.  The bound is uniform in the ground
function and volume, with the limiting convention of
Section~\ref{set:section}.
\end{proposition}

The proof constructs a common finite measure on the remaining particles
that has a density close to one relative to each of $P_v$ and $P_w$,
in a fourth-moment sense.  To construct it, we first represent deletion
by a distinguished label in the path law of Section~\ref{lp:section}.
The independent flags, exposed datum, and unexposed labels $\mathcal I$
are those defined there.  Two copies of a local path rearrangement
then give the same remaining configuration.  We estimate their
likelihoods before projecting back to $Y'$.

\subsection{The reference measures and count errors}

For a cell $v$ of either partition $\mathcal C_d$ or $\mathcal C_R$, put
\begin{equation}\label{palm:tag-count}
 T_v=f_0D|v|,\qquad h_{v,i}=\1_{\{X_i\in v\}},\qquad
 T_v^{\mathrm{rel}}=T_v^{-1}\sum_{i:\,i\text{ is flagged}}h_{v,i}.
\end{equation}
Here $X_i=\omega_i(0)$ is the midpoint of the path of label $i$.
Define a finite measure $P_{\mathrm{aug},v}$ on a path sample together
with a distinguished flagged label by
\begin{equation}\label{palm:augmented-reference}
 \int F\,\dd P_{\mathrm{aug},v}
 =\E\sum_{i:\,i\text{ is flagged}}\frac{h_{v,i}}{T_v}F(\omega,i).
\end{equation}
The expectation includes the flags.  Additional independent randomizations
introduced below are understood as probability kernels on this space.
Delete the distinguished midpoint and put the other $N-1$ midpoints in
an independently chosen uniform order.  The following identity shows
that this projection realizes the analytic measure $P_v$ above.

\begin{lemma}[Density of the smoothed deletion measure]
\label{palm:reference}
The projection of $P_{\mathrm{aug},v}$ just described is $P_v$.
The mass of $P_v$ equals $\E T_v^{\mathrm{rel}}$.  Its average over
either partition is one.
\end{lemma}

\begin{proof}
The midpoint law is $\Phi^2\,\dd X$.  Flagging a particular label has
probability $f_0$, independently of this law.  Symmetry of $\Phi^2$ and
the uniform ordering make all $N$ summands in
\eqref{palm:augmented-reference} identical after projection.  Their
coefficient is $Nf_0/T_v=V/|v|$.  This identifies the projection with
$P_v$ in \eqref{palm:rms-density}.  The last assertion follows by summing the
cell indicators, since $N=DV$.
\end{proof}

We always average a cell-indexed estimate by the number of cells in its
partition.  For adjacent ordered pairs of $d$-cells the normalization is
also the number of cells; changing this to the number of ordered edges
changes only a fixed constant.

\begin{lemma}[Fluctuations of the number of flagged particles]
\label{palm:counts}
The count estimates of Proposition~\ref{en:energy-outputs} imply
\begin{align}
 \frac1{|\mathcal C_R|}\sum_{v\in\mathcal C_R}
       \E|T_v^{\mathrm{rel}}-1|^2
 &\le C\bigl(D^{-1-w/6}+D^{-1+x}R^{-3}\bigr),
       \label{palm:R-count-moment}\\
 \frac1{|\mathcal C_d|}\sum_{v\in\mathcal C_d}
       \E|T_v^{\mathrm{rel}}-1|^4
 &\le C_j\bigl(D^{-1-u/6}+D^{-2+2x}\bigr).
       \label{palm:d-count-moment}
\end{align}
For each fixed integer $q$, products of at most $q$ relative counts
can be bounded by $C_jD^{qx}$ at an error smaller than every prescribed
negative power of $D$.  This assertion holds in any of the averaged
estimates involving a fixed number of cells used below.
\end{lemma}

\begin{proof}
Conditional on the configuration, the flagged count $M_v$ in a cell is
binomial with parameters $n_v$ and $f_0$.  In particular,
\[
 \E[(M_v-f_0n_v)^2\mid X]\le f_0n_v,
 \qquad
 \E[(M_v-f_0n_v)^4\mid X]
       \le C\bigl(f_0n_v+(f_0n_v)^2\bigr).
\]
Write
\[
 T_v^{\mathrm{rel}}-1
 =\frac{M_v-f_0n_v}{f_0D|v|}
        +\left(\frac{n_v}{D|v|}-1\right).
\]
The quadratic and quartic inequalities for this sum give the displayed
bounds from the count moments of Proposition~\ref{en:energy-outputs}.  The binomial contribution
to the fourth moment is at most
$C_j(D^{-2+2x}+D^{-3+3x})$; the second term is smaller for $x<1$.
The required first and second moments of $n_v/(D|v|)$ follow either
from these count estimates or from Proposition~\ref{set:count-mgf}.

For the last assertion, use the event that every unit cell meeting the
indicated cells has count at most $CD$, where $C$ is a sufficiently large
fixed constant.  On that event $T_v^{\mathrm{rel}}\le C_j/f_0$.
Proposition~\ref{set:count-mgf} bounds the complement by a polynomial in $D$
times $e^{-cD}$, also for an $R$-cell since it meets $O(R^3)$ unit cells.
Hard packing bounds all counts in these bounded regions by a fixed
power of $D$.  Consequently every fixed product on the exceptional
event has expectation smaller than any prescribed power of $D^{-1}$.
\end{proof}

We use two omissions.  A distinguished label is retained only if it
belongs to $\mathcal I$.  A pair of neighboring $d$-cells is compared only when
all sites in a fixed neighborhood of the two cells are good in the
sense of Definition~\ref{lp:certificate}.  The neighborhood includes
every possible tag mean and every path used in the construction below.
It has radius $C b$, with $C$ fixed.  For two $R$-cells we use their
enlargements by $CR$ instead.  Write $\mathcal G_{v,w}$ for the resulting good-neighborhood event.
These are events of the exposed datum $\mathcal D$;
we do not add tests on unexposed paths to that datum.

\begin{lemma}[Cost of omissions]
\label{palm:omissions}
For every indicated cell $v$, the mass in
$P_{\mathrm{aug},v}$ of tags outside $\mathcal I$ is at most
$C_{x,j}D^{-1-1/8}$, uniformly in $v$.  The mass lost by a failure of
the required good neighborhood is at most
$C_{x,j}D^{-1-1/8+x}$ for a $d$-cell and at most
$C_xR^3D^{-1-1/8+x}$ for an $R$-cell.
Multiplying an omission by a fixed product of relative counts costs
at most a further factor $C_jD^{qx}$, with $q$ the number of factors.
The same convention applies to a fixed power of the omitted relative
count itself.
\end{lemma}

\begin{proof}
Apart from flagging, failure of membership in $\mathcal I$ means failure of an
endpoint separation or an individual motion test.  Flag independence
cancels the factor $f_0$ in $T_v$.  A tag in $v$ satisfying the motion
cutoffs has both endpoints in a bounded enlargement of $v$.
Equation~\eqref{lp:close-pairs} bounds the expected number of labels there
having a second endpoint within distance $s$ by $CD^2s^3$ per unit
volume.  Division by $D$ gives $CDs^3=CD^{-7/5}$.
The visitor and free-cutoff estimates in the proof of
Proposition~\ref{lp:good-probability} give errors of arbitrarily high fixed
power for the remaining motion failures.  Summing over the unit cells
meeting an $R$-cell and dividing by its volume leaves the same bound.
Both estimates are stronger than the stated $D^{-1-1/8}$ bound.

For a good-neighborhood failure use Proposition~\ref{lp:good-probability}
and a union bound.  There are $O_{x,j}(1)$ tests for a local comparison
and $O_x(R^3)$ tests at scale $R$.  On the count event in
Lemma~\ref{palm:counts}, the total reference weight is bounded by
$C_j/f_0$.  Its complement is negligible.  This proves the second
claim and the assertion about products.
If $M_v$ denotes the omitted relative count, then
$0\le M_v\le C_j/f_0$ on the same count event and
$M_v^q\le(C_j/f_0)^{q-1}M_v$.  This proves the final convention.
\end{proof}

\subsection{A coupling with the same remaining midpoint configuration}

Fix disjoint cells $v,w$.  Select distinct flagged tags $i,j$ with
weight
\begin{equation}\label{palm:two-tag-weight}
             \frac{h_{v,i}h_{w,j}}{T_vT_w}.
\end{equation}
Retain only tags in $\mathcal I$ and the specified good-data event.  For a tag
use the unit cell containing its bridge mean as its site.  Confinement
places that site within distance $Cb$ of its indicated cell.
A route and its subsequent label choices depend on the data and the
tags, but not on unexposed path values.  Write the resulting list as
\begin{equation}\label{palm:label-list}
                 i_0=i,i_1,\ldots,i_m=j.
\end{equation}
All labels are distinct; consecutive bridge means have bounded
distance.  Interior labels are drawn from the supplies at the route
sites, with the tags excluded.  Conditional on the data, choices at
different sites are independent.  Their maximal point probability
is $C/(Df_0)$.  No interior choice is needed at an endpoint site.

\begin{lemma}[Two-copy midpoint transport]
\label{palm:transport}
For adjacent $d$-cells on the retained data event one can choose
\eqref{palm:label-list} of length $m\le C_x$, with
$C_x=O(x^{-1/2})$, and construct two path copies with these properties:
\begin{enumerate}
\item All unselected paths are the same as in the original sample.
The left path of $i$ and the right path of $j$ are also unchanged.
\item Both copies satisfy every path constraint defining the
conditional law at the given datum.
\item For $1\le k\le m$, the left midpoint of $i_k$ equals the right
midpoint of $i_{k-1}$.  Thus, after deleting $i$ on the left and $j$
on the right, the midpoint configurations agree as unordered lists.
\item The construction is a probability kernel conditional on the
original sample and the route choices.
\end{enumerate}
Consequently, weighting by \eqref{palm:two-tag-weight}, averaging all
choices, and ordering the common remaining midpoint list uniformly
defines a single finite measure $\nu_{v,w}$ on the remaining particles.
\end{lemma}

\begin{proof}
Join the tag mean-sites by a coordinate route of nearest-neighbor sites,
erase loops, and choose one interior label at each interior site.
The tag means lie within distance $Cb$ of adjacent cells, so there are
$O(b)=O(x^{-1/2})$ entries.  If the sites coincide the list consists
just of the two tags.  The good-neighborhood event supplies all labels.
The interior supplies have size at least $cf_0D$, so excluding the
bounded number of tags preserves the claimed point bound.

For $k=1,\ldots,m$, sample a common midpoint uniformly in a unit cell
whose distance from both corresponding means is bounded.  Fix that cell
as a measurable function of the exposed endpoints alone, before sampling
any latent original or proposed paths.  Conditional
on that point, propose independently the left path of $i_k$ and the
right path of $i_{k-1}$ as two-leg free bridges.  Condition this pair
on its individual cutoffs and on avoidance, in each copy separately,
of paths retained in that copy and paths already inserted in that copy.
Paths to be inserted at later stages are not checked at this stage.
Proposition~\ref{lp:insertion} gives success probability at least $1/2$ for
every such proposal, including every allowed preceding history.
Thus each stage is a probability kernel.  Every pair of paths in a
finished copy is checked when the later of them is inserted.  The
midpoint identities and equality of the two remaining lists follow
directly from the cyclic shift of labels in \eqref{palm:label-list}.
\end{proof}

We next calculate the density of the left output.  This calculation
also applies when selected outside obstacles and their checks have
been deleted; this flexibility will allow independent centering of
several route constructions.

Fix the data, $i$, the route labels, and the paths outside
$S=\{i_1,\ldots,i_m\}$.  Denote these outside paths by $\xi$.
For $o=(o_l)_{l\in S}$ define $a_S(o;\xi)$ to be the product of the
individual cutoff indicators and all avoidance indicators involving
at least one label in $S$.  The conditional reference law is
\begin{equation}\label{palm:conditional-reference}
 \lambda_S(\dd o\mid\xi)
  = Z_S(\xi)^{-1}a_S(o;\xi)\prod_{l\in S}B_l(\dd o_l),
 \quad
 Z_S(\xi)=\int a_S(o;\xi)\prod_{l\in S}B_l(\dd o_l)>0.
\end{equation}
Checks wholly outside $S$ are irrelevant to this conditional formula.
Let $K_S(o;\xi,\dd z)$ be the marginal probability kernel of the left
changed paths in Lemma~\ref{palm:transport}.  Its density relative to
$\prod_{l\in S}B_l$ is denoted by $k_S(o;\xi,z)$.
The unchanged factor $h_{v,i}/T_v$ will remain outside this calculation.

\begin{lemma}[Cancellation of the conditional normalizer]
\label{palm:likelihood}
The left-output likelihood, with the weight $h_{w,j}$ included but
the divisor $T_w$ omitted, is, on admissible reference values,
\begin{equation}\label{palm:likelihood-formula}
 L_S(z;\xi)
  =\int a_S(o;\xi)h_{w,j}(o)
          k_S(o;\xi,z)\prod_{l\in S}B_l(\dd o_l).
\end{equation}
It satisfies
\begin{equation}\label{palm:centered-identity}
 F_S=L_S-h_{w,j},\qquad
       \E_\lambda[F_S\mid\omega_{S^c}]=0.
\end{equation}
These assertions hold after removing any specified collection of
outside checks, provided the same removals are used in the reference,
the original-path integral, and the insertion kernels.  Moreover,
\begin{equation}\label{palm:local-likelihood-sup}
                    0\le L_S\le C_xD^{C\sqrt x}.
\end{equation}
\end{lemma}

\begin{proof}
Before dividing by the reference density, the output density is the
right side of \eqref{palm:likelihood-formula} divided by $Z_S$.
The kernel produces only admissible left paths; on that support the
reference density is $1/Z_S$.  These factors cancel.  Integrating the
resulting likelihood against \eqref{palm:conditional-reference} and
using that $K_S$ has total mass one gives
\[
 \int L_S\,\dd\lambda_S
   =Z_S^{-1}\int a_S(o;\xi)h_{w,j}(o)\prod_{l\in S}B_l(\dd o_l)
   =\int h_{w,j}\,\dd\lambda_S.
\]
This proves the conditional centering, also with deleted checks.

For clarity, the density used here can be written directly from the
pair proposals.  If $C_k$ is the chosen unit cell and $g_l$ the free
midpoint density for label $l$, then the proposal for the $k$th pair,
viewed from its left coordinate $z_{i_k}$, is
\begin{equation}\label{palm:pair-proposal-disintegration}
 \frac{\1_{C_k}(z_{i_k}(0))}{g_{i_k}(z_{i_k}(0))}
       B_{i_k}(\dd z_{i_k})\,
       B_{i_{k-1}}^{\,z_{i_k}(0)}(\dd z'_{i_{k-1}}).
\end{equation}
Here $B_l^y$ is the two-leg bridge pinned at midpoint $y$.
Its restriction to the individual cutoffs, together with the
compatibility checks, is divided by a success probability at least
$1/2$.  The first density factor in
\eqref{palm:pair-proposal-disintegration} is bounded by $CD^{Cx}$
on $C_k$, by Corollary~\ref{lp:trial-density}.  Integrating the auxiliary
right paths successively, bounding each normalizer by $2$, and using
$m=O(x^{-1/2})$ gives \eqref{palm:local-likelihood-sup}.
The original integral has mass at most one.  The same argument works
with checks deleted because their success probabilities only improve.
\end{proof}

\subsection{The centered fourth moment}

The likelihood $L_S$ is centered around the second-tag indicator
$h_{w,j}$, rather than around one.  For a fixed first tag $i$ with
$h_{v,i}=1$, average
over the route and interior labels, denoting this average by an overbar,
and then sum over candidate second tags: these are the labels in
$\mathcal I$ whose mean-sites lie in the required enlargement of $w$.
Confinement includes every retained tag with $h_{w,j}\ne0$ in this
data-dependent candidate set.
The exact decomposition is
\[
 \sum_{j\ne i}\frac{\overline{L_S}}{T_w}
 =\sum_{j\ne i}\frac{h_{w,j}}{T_w}
       +\sum_{j\ne i}\frac{\overline{F_S}}{T_w}.
\]
The first sum is the relative count in $w$ restricted to retained tags;
Lemmas~\ref{palm:counts} and~\ref{palm:omissions} control its difference
from one.  We now estimate the fourth moment of the second sum.
Expanding it produces four centered route likelihoods with a common
first tag.  If one of their changed label sets is disjoint from the
others and all its interactions with them are removed, its conditional
integral vanishes.  When all four changed sets are disjoint, a nonzero
term therefore needs at least two interactions between lists.  The
next lemma quantifies their joint cost, and also the one-interaction
cost needed when two lists share a label.

We specify the removed interactions for four possibly overlapping
route lists, all with this same first tag $i$.
Each list has a changed set $S_a$, $1\le a\le4$; these
sets exclude $i$.  Put $U=\bigcup_aS_a$ and condition on paths outside
$U$.  Join two lists when their changed sets intersect, and call the
connected components of this relation blocks.  A singleton block
contains just one list, whether or not that list has several labels.

For each unordered pair $e$ of different blocks introduce a switch
$\sigma_e\in\{0,1\}$.  Turning this switch off removes the following
checks everywhere in the conditional integral:
\begin{enumerate}
\item reference avoidance checks between the two blocks;
\item checks between original paths integrated in one likelihood and
outside reference paths in the other block;
\item checks, and the corresponding checks in insertion success
probabilities, between proposed paths in one likelihood and outside
reference paths in the other block.
\end{enumerate}
Individual cutoffs and all other checks are retained.  A newly proposed
right path with label $i$ belongs to the block generating that proposal.
Thus its checks against another block are switched, whereas checks of left insertions
against the unchanged left path of $i$ are always fixed-outside checks.
A right proposal is never checked against that left copy of $i$.
The right retained path of the second tag is an original path of its
own block.  This specifies every check even though the two path copies
retain different endpoint tags.

Let $F_a^{\sigma}$ be the centered likelihood constructed using these
switches.  Lemma~\ref{palm:likelihood} applies for every switch setting.
All formulas below hold on a fixed conditional fiber; constants are
uniform over its allowed outside paths and good datum.

\begin{lemma}[Two switched interactions]
\label{palm:two-switches}
Write $\beta_{lk}$ for the collision bound in
Lemma~\ref{lp:collision}, and set
\[
       \beta_* = C_xD^{-19/100}.
\]
Suppose either that the four changed sets are disjoint, or that they
have exactly one independent label identification.  After increasing $D$
the following bounds hold for their centered likelihoods, integrated
against the conditional reference on their union.  Finite differences
act on the entire admissibility numerator and its likelihood factors;
the original conditional-reference normalizer is held fixed.
\begin{enumerate}
\item A finite difference that switches one block-pair interaction
has absolute integral at most
$C_xD^{C\sqrt x}$ times the sum of the $\beta_{lk}$ for the finitely
many label pairs checked across that interaction.
In particular it is at most $C_xD^{C\sqrt x}\beta_*$.
\item If all four lists have disjoint changed sets, a mixed difference
in two distinct block-pair interactions has absolute integral at most
$C_xD^{C\sqrt x}$ times a finite sum of products
$\beta_{lk}\beta_{l'k'}$, with one checked pair from each interaction.
\end{enumerate}
The bounds are valid if additional switches have already been fixed
on or off.  The number of summands is bounded by a constant depending
only on $x$.
\end{lemma}

\begin{proof}
Expand each centered factor as $L_S-h_{w,j}$ and use
\eqref{palm:likelihood-formula} and
\eqref{palm:pair-proposal-disintegration} for the likelihood factors.
This represents the conditional integral using the following variables:
reference output bridges on $U$; a separate collection of original
bridges for each likelihood; and the auxiliary right bridges in its
pair proposals.  Without avoidance checks these variables factor
between blocks.  Within one block the two members of a pair proposal
have a common midpoint, as they must.  Every indicator imposing an
individual cutoff is retained in the estimates below.

There are three kinds of factors: indicators, success-probability
reciprocals, and the densities in
\eqref{palm:pair-proposal-disintegration}.  The last factors do not
depend on switches.  When the changed sets are disjoint, combine each
such density with its free left-output bridge and its pinned auxiliary
bridge: this is exactly the probability measure of the common-uniform-
midpoint pair proposal.  We retain that measure when integrating a
collision.  In particular we do not replace it by the product of two
independent bridges; that product would not even dominate the pair law.
Thus the base measure within each singleton block is a product of plain
original bridges and independent pair proposals, or of plain output
bridges if its centered factor contributes the subtracted indicator.
The different blocks have independent base measures, conditional on the
fixed outside paths.  Every proposal cell was fixed using the exposed
endpoints, so its law has no remaining dependence on other blocks.

There are $O(x^{-1/2})$ pairs, including all four lists.  All success
reciprocals are at most two.  The inverse reference normalizer costs a
fixed factor $C_x$, by successive single-path insertions on the bounded
set $U$; keep that normalizer fixed when taking differences.  With one
identification, the nonsingleton block has repeated output coordinates.
We leave those coordinates and their pinned auxiliary variables together
until the singleton collision integrations have been carried out.
Afterwards the remaining integral is bounded by the marginal-density
argument proving \eqref{palm:local-likelihood-sup}, paying at most
$C_xD^{C\sqrt x}$.  This estimate never requires a joint density for
two paths with a common midpoint.  These observations also bound all
unused factors in a difference by the same displayed expense.

An indicator difference caused by a block-pair switch is supported
on the event that at least one of the removed avoidance checks fails.
Splitting this event by the union bound gives a collision between
one path from each of the two blocks.  The paths have distinct labels
and meet their individual cutoffs.  Integrating a free bridge or a
marginal of a common-uniform-midpoint proposal against a deterministic
allowed opponent costs at most $C_x\beta_{lk}$, by
Lemma~\ref{lp:collision}.  The same assertion holds with the roles
of the two paths exchanged.

We give the analogous verification for denominators.  At a fixed
insertion stage and preceding history its success probability is
\[
 p(\sigma)=\int Q(\dd z)\,c(z)\prod_e A_e(z)^{\sigma_e},
 \qquad p(\sigma)\ge\tfrac12.
\]
Here $Q$ is a fresh pair proposal, $c$ contains the individual cutoffs
and unswitched checks, and $A_e$ is the product of checks controlled
by switch $e$ at that stage.  The convention $A^0=1$ applies also
when $A=0$.  A first difference of $p$ is an integral with the factor
$A_e-1$; a mixed difference in $e,f$ has the factor
$(A_e-1)(A_f-1)$.  They therefore force one or both corresponding
collisions.  The reciprocal identity
\[
          p^{-1}-q^{-1}=\frac{q-p}{pq}
\]
and its difference in a second variable bound a first reciprocal
difference by a constant times $|\Delta_ep|$, and a second difference
by a constant times
\[
 |\Delta_e\Delta_fp|
 +\max_{\sigma}|\Delta_ep(\sigma)|\,
       \max_{\sigma}|\Delta_fp(\sigma)|.
\]
The product is represented using independent fresh proposals.  These
proposals belong to the block whose denominator is being tested.
After dropping unused compatibility indicators, they are ordinary
pair proposals with their own cutoff indicators.  They introduce only
a fixed number of additional density factors.

Apply the finite-product difference rule.  A first difference is
thus a finite sum of integrals with one forced cross-block collision.
To estimate it, integrate a singleton block against its fixed
opponents whenever there is an overlapping block.  In the application
with one identification every block pair has a singleton side unless
both blocks are singletons, which is easier.  The remaining factors
are bounded as above.  This proves the first assertion in precisely
the overlapping case needed below; with four disjoint lists either
side can be integrated.

For the second assertion the two different block-pair edges form
either a two-edge tree or two disjoint edges.  In the tree case fix
all variables in the central block and integrate the two leaf blocks
in succession.  Each leaf participates in just one of the two forced
collisions, so its integration costs the appropriate $\beta$ bound
uniformly in the central block.  No independence of collisions inside
the central block is used.  If a forced collision comes from a
denominator test, its fresh proposal belongs to its generating block;
the same leaf integration applies.  The two disjoint edges are
estimated separately.  This proves the product bound.

Extra differences in other switches, when they occur, are bounded
by a finite sum of evaluations of the same two-difference expression.
This does not increase the number of forced collisions that must be
estimated.  Finally $\beta_{lk}\le b_D(1+2/s)$, which is bounded by
$C_xD^{-19/100}$ for sufficiently large $D$.  This proves the stated
worst bound.
\end{proof}

\begin{lemma}[The fourth moment of the averaged centered likelihood]
\label{palm:fourth}
For adjacent $d$-cells $v,w$, let an overbar denote the average over
the route and interior label choices at fixed tags and datum.  Restrict
the expectation to $\mathcal G_{v,w}$ and retained first tags.  Equivalently,
set the complete centered summand to zero outside this event.  Then
\begin{equation}\label{palm:fourth-centered}
 \E\left[\1_{\mathcal G_{v,w}}
       \sum_{i\in\mathcal I}\frac{h_{v,i}}{T_v}
    \left|\sum_{j\ne i}\frac{\overline{F_S}}{T_w}\right|^4\right]
          \le C_{x,j}D^{-1-1/20},
\end{equation}
provided $x$ is sufficiently small.  The candidate second tags are
the labels in $\mathcal I$ whose mean-sites lie in the required enlargement
of $w$.
\end{lemma}

\begin{proof}
Expand the fourth power into four independently randomized routes
and label lists, conditionally on the datum and the common tag $i$.
The total tag counting weights are bounded by $C_{x,j}$ on good data:
there are at most $C_xDf_0$ candidates in the bounded enlargement,
and $T_v,T_w$ are fixed multiples of $Df_0$.  Each normalized tag
weight, and each interior label probability, is at most
$C_{x,j}/(Df_0)$.

First split the tag summations by the mean-sites of the second tags
and the resulting route site patterns, which have only $O_x(1)$
possibilities.  Retain the original summation weights in each subsum;
there is no renormalization by the mass of a site pattern.  For every
fixed tag choice in such a subsum, each tag counting weight and each
interior-choice probability is bounded by $C_{x,j}/(Df_0)$.  The set of
candidate labels for any slot has at most $C_xDf_0$ elements.  In bounding
a positive sum we may discard within-list exclusions and dominate its
joint weight by the product of these point bounds on the respective
candidate sets.  This domination requires no claim that conditioning on
an arbitrary previously revealed tag preserves the stated point-mass bound.

Two or more independent identifications between the four changed
lists therefore cost at most $C_{x,j}(Df_0)^{-2}$.  An identification means
one equality needed to merge previously separate label slots; a triple
equality counts as two.  Indeed an equality pattern with $q$ independent
identifications leaves $q$ fewer freely chosen labels.  Summing the
product point bounds over that pattern gives at most
$C_{x,j}(Df_0)^{-q}$, since every free label has at most $C_xDf_0$ possible
values.  There are only $O_x(1)$ slot patterns.  A label cannot occur
twice in one list, and none of the changed lists contains $i$.
The supremum bound
\eqref{palm:local-likelihood-sup} therefore bounds all contributions
with at least two identifications by
\begin{equation}\label{palm:many-matches}
                  C_{x,j}D^{-2+2x+C\sqrt x}.
\end{equation}

There remain the cases of no identification and exactly one.
Condition on paths outside $U$ and write the conditional expectation
as $Z_U^{-1}$ times an integral against the free product bridges on
$U$ with its admissibility numerator.  Keep $Z_U$ fixed.  Regard
this numerator, including the four $F_a^\sigma$, as a function $G$ of
the block-pair switches.  The elementary expansion at the all-off
vertex of the Boolean cube is
\begin{equation}\label{palm:switch-expansion}
          G(\boldsymbol1)=\sum_{E}\Delta_EG(\boldsymbol0).
\end{equation}
In a term with an isolated singleton block, the integrand factors
with respect to that block and the remaining blocks.  Integrating
the isolated block gives zero by
\eqref{palm:centered-identity}, multiplied by its unnormalized
reference density.  This is a statement about each term of the
mixed difference: its absent edges are off in every evaluation.

With exactly one identification there are three blocks, two of them
singletons.  A surviving term in \eqref{palm:switch-expansion} must
therefore contain at least one edge.  Lemma~\ref{palm:two-switches}
and the identification cost bound its total by
\begin{equation}\label{palm:one-match}
                     C_{x,j}D^{-1+x-19/100+C\sqrt x}.
\end{equation}
With no identification all four blocks are singletons.  A surviving
term must contain at least two different edges.  Choose two of them,
apply the two-difference bound, and estimate the remaining differences
by the finite sum described in that lemma.

One of the two collision factors can be averaged over a free label
choice.  There is always such a label: every cross-block check has a
label of a changed list on at least one side, whereas $i$ belongs to
none of those lists.  The flagged-endpoint bound in
Definition~\ref{lp:certificate}, summed over dyadic distance annuli,
gives, uniformly in the other endpoint,
\begin{equation}\label{palm:averaged-beta}
 \frac{C_{x,j}}{Df_0}\sum_{k\text{ in a candidate list}}
       \frac1{s\vee d(Y_k,y)}
 \le C_{x,j}\left(1+\frac{D^{1/100}}{Df_0s}\right)\le C_{x,j}.
\end{equation}
The same holds at the other endpoint time.  The candidate list is
contained in a fixed bounded neighborhood; the annulus estimate uses
$\#B_h\le C(Df_0h^3+D^{1/100})$.  The last inequality follows from
$x<19/100$.  Thus the averaged collision factor is at most $C_{x,j}b_D$.
We bound the other factor by $\beta_*$.  Excluding additional label
coincidences can only decrease these positive upper bounds, so their
restriction need not be retained during this summation.  The no-match
contribution is consequently
\begin{equation}\label{palm:no-match}
                       C_{x,j}D^{-1-19/100+C\sqrt x}(\log D)^C.
\end{equation}
The factor $h_{v,i}$ depends only on the fixed outside path $i$ and
can be retained throughout the cancellations.  It is bounded by one
when estimating the errors.  Summing the first tags has total weight
at most $C_{x,j}$ on good data.  Combining \eqref{palm:many-matches},
\eqref{palm:one-match}, and \eqref{palm:no-match} proves
\eqref{palm:fourth-centered} when $x$ is small enough that all power
losses leave a saving of $1/20$.
\end{proof}

\subsection{Projection to the remaining particles}

The measure $\nu_{v,w}$ has two lifts: one to
$P_{\mathrm{aug},v}$, observing the left copy and tag $i$, and one
to $P_{\mathrm{aug},w}$, observing the right copy and tag $j$.
Write their respective densities as $\ell_v$ and $\ell_w$.
On the retained part of the left reference,
\begin{equation}\label{palm:lift-density}
       \ell_v=\sum_{j\ne i}T_w^{-1}\overline{L_S}.
\end{equation}
It is zero elsewhere.  The sum of the corresponding centered
subtractions is the relative count in $w$ restricted to tags in $\mathcal I$.
Indeed, confinement includes every such tag with nonzero $h_{w,j}$
in the candidate mean range, and $h_{v,i}h_{w,i}=0$ because the cells
are disjoint.

\begin{proposition}[Fourth-moment comparison before projection]
\label{palm:augmented-comparison}
For sufficiently small $x>0$, depending on the previously chosen
exponent $u$, the measures above satisfy
\begin{equation}\label{palm:augmented-fourth}
 \frac1{|\mathcal C_d|}\sum_{v\sim w}
     \left\{\int|\ell_v-1|^4\,\dd P_{\mathrm{aug},v}
          +\int|\ell_w-1|^4\,\dd P_{\mathrm{aug},w}\right\}
       \le C_jD^{-1-u/8}.
\end{equation}
The summation can be restricted to adjacent cells in the same $R$-cell.
The bound is uniform in volume and in the nonnegative ground function,
with the limiting convention of Section~\ref{set:section}.
\end{proposition}

\begin{proof}
Let $M_w$ be the omitted relative count in $w$.  On the retained left
reference, \eqref{palm:lift-density} gives
\[
 \ell_v-1=\sum_{j\ne i}T_w^{-1}\overline{F_S}
                  +(T_w^{\mathrm{rel}}-1)-M_w.
\]
The inequality $|a+b+c|^4\le27(|a|^4+|b|^4+|c|^4)$ separates the
three terms.  Lemma~\ref{palm:fourth} treats the first.  For the second,
the count-tail bound gives
\[
 \E\bigl[T_v^{\mathrm{rel}}|T_w^{\mathrm{rel}}-1|^4\bigr]
 \le C_jD^x\E|T_w^{\mathrm{rel}}-1|^4+O(D^{-A})
\]
for any fixed $A$, after averaging cells.  Neighboring pairs have
uniformly bounded degrees, so Lemma~\ref{palm:counts} applies.
For the third term use
$\E[T_v^{\mathrm{rel}}M_w^4]
 \le C_jD^{4x}\E M_w+O(D^{-A})$
and Lemma~\ref{palm:omissions}.
On omitted first tags or data the density is zero, so the fourth
error is exactly the omitted reference mass.  The same lemma bounds
it.  Choose $x$ so that, in addition to the requirements of
Lemma~\ref{palm:fourth},
\[
       x<u/24,\qquad 4x+u/8<1/8.
\]
The resulting powers are bounded by $D^{-1-u/8}$; the binomial
$D^{-2+2x}$ contribution is smaller.  This proves the left half of
\eqref{palm:augmented-fourth}.  For the right half, analyze the same
generated kernel with the observed and auxiliary copies exchanged,
retaining its original insertion order.  Every preceding density and
switch estimate holds for any fixed insertion order: only the specified
previous-path checks enter its success probabilities.  Thus this argument
uses no symmetry of the success normalizers.  Interchanging the two tag
roles gives the other half.
\end{proof}

\begin{proof}[Proof of Proposition~\ref{palm:local-comparison}]
Let $q$ be the Lebesgue density of the common projection $\nu_{v,w}$.
Conditional Jensen under the projection gives
\[
 \int\left|\frac q{s_v^2}-1\right|^4s_v^2\,\dd Y'
       \le\int|\ell_v-1|^4\,\dd P_{\mathrm{aug},v},
\]
with the analogous estimate for $w$.  The measure $\nu_{v,w}$ is
absolutely continuous with respect to both references, so $q=0$
where either corresponding density vanishes.

For completeness, for $a,b\ge0$ and $q\ge0$ with this zero-set
convention there is an absolute constant $C$ such that
\begin{equation}\label{palm:scalar-fourth-transfer}
 \frac{(a-b)^4}{a^2+b^2}
 \le C\left\{a^2\left|\frac q{a^2}-1\right|^4
            +b^2\left|\frac q{b^2}-1\right|^4\right\}.
\end{equation}
For $a,b>0$, put $c=\sqrt q$.  The inequalities
$|a-b|^4\le8(|a-c|^4+|b-c|^4)$ and
$|1-\sqrt z|\le|1-z|$ for $z\ge0$ give
\[
 \frac{|a-c|^4}{a^2+b^2}
 \le a^2\left|1-\frac c a\right|^4
 \le a^2\left|1-\frac q{a^2}\right|^4,
\]
with the same estimate for $b$.  This proves
\eqref{palm:scalar-fourth-transfer} with $C=8$.  If one of $a,b$
vanishes, absolute continuity forces $q=0$, and the inequality follows
directly; if both vanish, both sides are defined to be zero.

Apply \eqref{palm:scalar-fourth-transfer} to $a=s_v$, $b=s_w$ and sum
the Jensen bounds.  This proves \eqref{palm:local-rms-comparison}.
\end{proof}
